\input format.tex

\usepackage{graphicx}
\graphicspath{{cores/}}

\usepackage{environ}
\usepackage{colortbl,array,booktabs}
\usepackage{tabularx}

\colorlet{TablaBordeSuperior}{topcolor}
\colorlet{TablaBordeInferior}{topcolor}
\colorlet{TablaCentroSuperior}{blue!1}
\colorlet{TablaCentroInferior}{blue!20}
\colorlet{FuenteCabeceraTabla}{white}

\newcolumntype{M}[1]{>{\centering\arraybackslash}m{#1}}

\tcbset{rtab/.style={
freelance,
frame code={
 \path[top color=topcolor,bottom color=topcolor]
   ([yshift=-#1*(\baselineskip+2pt)]interior.north west) --
   ([yshift=-#1*(\baselineskip+2pt)]interior.north east) {[sharp corners]--
    ([yshift=3pt]interior.north east) --
    ([yshift=3pt]interior.north west)} -- cycle;

  },
interior code={},
 }
}

\newcommand\fuentecabecera[1]{\textcolor{black}{\textbf{#1}}}

\NewEnviron{RCtable}[1]
  {%
    \begin{tcolorbox}[#1]\BODY\end{tcolorbox}%
  }



\begin{document}

\vspace*{8mm}
%% 各章节
\setlength{\arrayrulewidth}{1pt}
\fontsize{9.3pt}{11pt}\selectfont
\color{gray2}

\begin{center}
{\noindent\bf\sanhao 检测结果}
\end{center}

\vspace*{4mm}
\noindent
{\noindent\bf\xiaosihao 1、肠道菌群概况}
\vspace*{1.5mm}

%\begin{RCtable}{
%  rtab=1,
%  tabularx*={\arrayrulecolor{gray2}}%
%    {M{0.3\linewidth} M{0.31\linewidth}M{0.34\linewidth}}
%}
%\fuentecabecera{测试项目} & \fuentecabecera{健康状态提示} & \fuentecabecera{检测结果评价} \\[3pt]
%\bf{检测项目} & \bf{健康状态提示} & \bf{检测结果评价} \\[3pt]
%\end{RCtable}

\begin{tctabularx}{tabularx={m{4.8cm}<{\centering}m{5.2cm}<{\centering}m{4.55cm}<{\centering}}}
&&
\\[-6pt]
\cellcolor{topcolor} \raisebox{2.614pt}{\color{black00}\bf 检测项目} &
\cellcolor{topcolor} \raisebox{2.614pt}{\color{black00}\bf 健康状态提示} &
\cellcolor{topcolor} \raisebox{2.614pt}{\color{black00}\bf 检测结果评价}
\end{tctabularx}

\vspace*{-4.25mm}
\fontsize{8.8pt}{11pt}\selectfont
\arrayrulecolor{gray2}
\begin{longtable}{m{4.8cm}m{5.2cm}<{\centering}m{0cm}@{}m{4.61cm}<{\centering}}
\hline
\parbox[c]{\hsize}{\vskip7pt {\lantxh 肠道菌群多样性\\用于衡量肠道菌群物种丰富度及稳定性的指标} \vskip7pt} & \parbox[c]{\hsize}{\vskip7pt\centerline{\raisebox{-1.5ex}{\includegraphics[width=5cm,height=0.55cm]{bar.pdf}}}\vskip7pt}  &
\hspace*{-1.51cm}\raisebox{-0.45ex}{\includegraphics{cry.pdf}}
 & \begin{minipage}{4.60cm}\begin{center}{{\color{red}\lantxh 低{\\ \bahao 肠道菌群组成较为单一，菌群失调风险高。}} }\end{center} \end{minipage} \\
\hline
\parbox[c]{\hsize}{\vskip7pt {\lantxh 有益菌\\产生有益物质，抑制病原菌的生长，帮助预防疾病} \vskip7pt} & \parbox[c]{\hsize}{\vskip7pt\centerline{\raisebox{-1.5ex}{\includegraphics[width=5cm,height=0.55cm]{bar.pdf}}}\vskip7pt}  &
\hspace*{-1.51cm}\raisebox{-0.45ex}{\includegraphics{cry.pdf}}
 & \begin{minipage}{4.60cm}\begin{center}{{\color{red}\lantxh 低{\\ \bahao 有益物质产出不足，不利于肠道健康}} }\end{center} \end{minipage} \\
\hline
\parbox[c]{\hsize}{\vskip7pt {\lantxh 有害菌\\产生多种毒素，危害肠道健康} \vskip7pt} & \parbox[c]{\hsize}{\vskip7pt\centerline{\raisebox{-1.5ex}{\includegraphics[width=5cm,height=0.55cm]{bar.pdf}}}\vskip7pt}  &
\hspace*{-4.83cm}\raisebox{-0.45ex}{\includegraphics{smile.pdf}}
 & \begin{minipage}{4.60cm}\begin{center}{{\lantxh 低{\\ \bahao 有害菌含量低，有益于肠道健康}} }\end{center} \end{minipage} \\
\hline
\end{longtable}

\vspace*{2mm}
\noindent
{\noindent\bf\xiaosihao 2、常见肠道菌含量}
\vspace*{1.5mm}

\begin{tctabularx}{tabularx={m{4.8cm}<{\centering}m{5.2cm}<{\centering}m{4.55cm}<{\centering}}}
&&
\\[-6pt]
\cellcolor{topcolor} \raisebox{2.614pt}{\color{black00}\bf 检测项目} &
\cellcolor{topcolor} \raisebox{2.614pt}{\color{black00}\bf 健康状态提示} &
\cellcolor{topcolor} \raisebox{2.614pt}{\color{black00}\bf 检测结果评价}
\end{tctabularx}

\vspace*{-4.25mm}
\fontsize{8.8pt}{11pt}\selectfont
\arrayrulecolor{gray2}
\begin{longtable}{m{4.8cm}m{5.2cm}<{\centering}m{0cm}@{}m{4.61cm}<{\centering}}
\hline
\parbox[c]{\hsize}{\vskip7pt {\lantxh 拟杆菌属\\常见肠道菌，多数与人体互利共生，少数菌种可能会导致菌血症等} \vskip7pt} & \parbox[c]{\hsize}{\vskip7pt\centerline{\raisebox{-1.5ex}{\includegraphics[width=5cm,height=0.55cm]{bar.pdf}}}\vskip7pt}  &
\hspace*{-1.51cm}\raisebox{-0.45ex}{\includegraphics{cry.pdf}}
 & \begin{minipage}{4.60cm}\begin{center}{{\color{red}\lantxh 低{\\ \bahao 不利于肠道菌群平衡}} }\end{center} \end{minipage} \\
\hline
\parbox[c]{\hsize}{\vskip7pt {\lantxh 布劳特氏菌属\\发酵多种植物多糖产生乙酸盐，促进肠道健康} \vskip7pt} & \parbox[c]{\hsize}{\vskip7pt\centerline{\raisebox{-1.5ex}{\includegraphics[width=5cm,height=0.55cm]{bar.pdf}}}\vskip7pt}  &
\hspace*{-1.51cm}\raisebox{-0.45ex}{\includegraphics{cry.pdf}}
 & \begin{minipage}{4.60cm}\begin{center}{{\color{red}\lantxh 低{\\ \bahao 不利于产生有益物质及肠道健康}} }\end{center} \end{minipage} \\
\hline
\parbox[c]{\hsize}{\vskip7pt {\lantxh 栖粪杆菌属\\发酵纤维素产生丁酸等有益物质，抑制肠道炎症，促进肠道健康} \vskip7pt} & \parbox[c]{\hsize}{\vskip7pt\centerline{\raisebox{-1.5ex}{\includegraphics[width=5cm,height=0.55cm]{bar.pdf}}}\vskip7pt}  &
\hspace*{-1.51cm}\raisebox{-0.45ex}{\includegraphics{cry.pdf}}
 & \begin{minipage}{4.60cm}\begin{center}{{\color{red}\lantxh 低{\\ \bahao 不利于产生有益物质及抑制肠道炎症}} }\end{center} \end{minipage} \\
\hline
\parbox[c]{\hsize}{\vskip7pt {\lantxh 粪球菌属\\发酵多糖，能够产生丁酸、乙酸等有益物质，调节肠道环境} \vskip7pt} & \parbox[c]{\hsize}{\vskip7pt\centerline{\raisebox{-1.5ex}{\includegraphics[width=5cm,height=0.55cm]{bar.pdf}}}\vskip7pt}  &
\hspace*{-1.51cm}\raisebox{-0.45ex}{\includegraphics{cry.pdf}}
 & \begin{minipage}{4.60cm}\begin{center}{{\color{red}\lantxh 低{\\ \bahao 不利于产生有益物质及调节肠道环境}} }\end{center} \end{minipage} \\
\hline
\parbox[c]{\hsize}{\vskip7pt {\lantxh 瘤胃球菌属\\帮助降解纤维素等人体不能消化的多糖，在肠道中广泛分布} \vskip7pt} & \parbox[c]{\hsize}{\vskip7pt\centerline{\raisebox{-1.5ex}{\includegraphics[width=5cm,height=0.55cm]{bar.pdf}}}\vskip7pt}  &
\hspace*{-1.51cm}\raisebox{-0.45ex}{\includegraphics{cry.pdf}}
 & \begin{minipage}{4.60cm}\begin{center}{{\color{red}\lantxh 低{\\ \bahao 不利于降解纤维素等多糖}} }\end{center} \end{minipage} \\
\hline
\parbox[c]{\hsize}{\vskip7pt {\lantxh 颤螺菌属\\帮助抗性淀粉和脂肪消化，保持正常体重，抑制肠道炎症} \vskip7pt} & \parbox[c]{\hsize}{\vskip7pt\centerline{\raisebox{-1.5ex}{\includegraphics[width=5cm,height=0.55cm]{bar.pdf}}}\vskip7pt}  &
\hspace*{-1.51cm}\raisebox{-0.45ex}{\includegraphics{cry.pdf}}
 & \begin{minipage}{4.60cm}\begin{center}{{\color{red}\lantxh 低{\\ \bahao 不利于保持正常体重及抑制肠道炎症}} }\end{center} \end{minipage} \\
\hline
\parbox[c]{\hsize}{\vskip7pt {\lantxh 副拟杆菌属\\帮助消化纤维素、抗性淀粉，保护肠道，抑制肠道炎症} \vskip7pt} & \parbox[c]{\hsize}{\vskip7pt\centerline{\raisebox{-1.5ex}{\includegraphics[width=5cm,height=0.55cm]{bar.pdf}}}\vskip7pt}  &
\hspace*{-1.51cm}\raisebox{-0.45ex}{\includegraphics{cry.pdf}}
 & \begin{minipage}{4.60cm}\begin{center}{{\color{red}\lantxh 低{\\ \bahao 不利于抑制肠道炎症及保护肠道}} }\end{center} \end{minipage} \\
\hline
\parbox[c]{\hsize}{\vskip7pt {\lantxh 毛螺菌属\\发酵多种糖类产生乙酸、甲酸等物质，能保护肠黏膜，抑制肠道炎症} \vskip7pt} & \parbox[c]{\hsize}{\vskip7pt\centerline{\raisebox{-1.5ex}{\includegraphics[width=5cm,height=0.55cm]{bar.pdf}}}\vskip7pt}  &
\hspace*{-1.51cm}\raisebox{-0.45ex}{\includegraphics{cry.pdf}}
 & \begin{minipage}{4.60cm}\begin{center}{{\color{red}\lantxh 低{\\ \bahao 不利于产生有益物质及抑制肠道炎症}} }\end{center} \end{minipage} \\
\hline
\parbox[c]{\hsize}{\vskip7pt {\lantxh 链球菌属\\多数为共生菌，少数在免疫力弱时可能引起化脓性炎症、肺炎等} \vskip7pt} & \parbox[c]{\hsize}{\vskip7pt\centerline{\raisebox{-1.5ex}{\includegraphics[width=5cm,height=0.55cm]{bar.pdf}}}\vskip7pt}  &
\hspace*{-1.51cm}\raisebox{-0.45ex}{\includegraphics{cry.pdf}}
 & \begin{minipage}{4.60cm}\begin{center}{{\color{red}\lantxh 低{\\ \bahao 不利于肠道菌群平衡}} }\end{center} \end{minipage} \\
\hline
\parbox[c]{\hsize}{\vskip7pt {\lantxh 普雷沃氏菌属\\肠道共生菌，偏素食人群中含量高，少数菌种与感染、牙周疾病有关} \vskip7pt} & \parbox[c]{\hsize}{\vskip7pt\centerline{\raisebox{-1.5ex}{\includegraphics[width=5cm,height=0.55cm]{bar.pdf}}}\vskip7pt}  &
\hspace*{-1.51cm}\raisebox{-0.45ex}{\includegraphics{cry.pdf}}
 & \begin{minipage}{4.60cm}\begin{center}{{\color{red}\lantxh 低{\\ \bahao 不利于肠道菌群平衡}} }\end{center} \end{minipage} \\
\hline
\parbox[c]{\hsize}{\vskip7pt {\lantxh 罗斯拜瑞氏菌属\\产生丁酸等有益物质，抑制肠道炎症，有利于肠道及人体健康} \vskip7pt} & \parbox[c]{\hsize}{\vskip7pt\centerline{\raisebox{-1.5ex}{\includegraphics[width=5cm,height=0.55cm]{bar.pdf}}}\vskip7pt}  &
\hspace*{-1.51cm}\raisebox{-0.45ex}{\includegraphics{cry.pdf}}
 & \begin{minipage}{4.60cm}\begin{center}{{\color{red}\lantxh 低{\\ \bahao 不利于产生有益物质及抑制肠道炎症}} }\end{center} \end{minipage} \\
\hline
\parbox[c]{\hsize}{\vskip7pt {\lantxh 双歧杆菌属\\有益菌，降解人体不能消化的多糖，产乳酸，调节免疫及肠道环境} \vskip7pt} & \parbox[c]{\hsize}{\vskip7pt\centerline{\raisebox{-1.5ex}{\includegraphics[width=5cm,height=0.55cm]{bar.pdf}}}\vskip7pt}  &
\hspace*{-1.51cm}\raisebox{-0.45ex}{\includegraphics{cry.pdf}}
 & \begin{minipage}{4.60cm}\begin{center}{{\color{red}\lantxh 低{\\ \bahao 不利于调节免疫及肠道环境}} }\end{center} \end{minipage} \\
\hline
\parbox[c]{\hsize}{\vskip7pt {\lantxh 阿克曼氏菌属\\降解粘蛋白、调节免疫，有利于肠黏膜完整性，保持正常体重} \vskip7pt} & \parbox[c]{\hsize}{\vskip7pt\centerline{\raisebox{-1.5ex}{\includegraphics[width=5cm,height=0.55cm]{bar.pdf}}}\vskip7pt}  &
\hspace*{-1.51cm}\raisebox{-0.45ex}{\includegraphics{cry.pdf}}
 & \begin{minipage}{4.60cm}\begin{center}{{\color{red}\lantxh 低{\\ \bahao 不利于调节免疫及保持正常体重}} }\end{center} \end{minipage} \\
\hline
\parbox[c]{\hsize}{\vskip7pt {\lantxh 柯林斯氏菌属\\发酵多种糖类产生氢气、乙醇等} \vskip7pt} & \parbox[c]{\hsize}{\vskip7pt\centerline{\raisebox{-1.5ex}{\includegraphics[width=5cm,height=0.55cm]{bar.pdf}}}\vskip7pt}  &
\hspace*{-1.51cm}\raisebox{-0.45ex}{\includegraphics{cry.pdf}}
 & \begin{minipage}{4.60cm}\begin{center}{{\color{red}\lantxh 低{\\ \bahao 不利于发酵糖类}} }\end{center} \end{minipage} \\
\hline
\parbox[c]{\hsize}{\vskip7pt {\lantxh 韦荣氏球菌属\\分解乳酸为乙酸和丙酸，调节肠道pH，极少数菌种可能引起炎症} \vskip7pt} & \parbox[c]{\hsize}{\vskip7pt\centerline{\raisebox{-1.5ex}{\includegraphics[width=5cm,height=0.55cm]{bar.pdf}}}\vskip7pt}  &
\hspace*{-1.51cm}\raisebox{-0.45ex}{\includegraphics{cry.pdf}}
 & \begin{minipage}{4.60cm}\begin{center}{{\color{red}\lantxh 低{\\ \bahao 不利于产生有益物质及调节肠道pH值}} }\end{center} \end{minipage} \\
\hline
\parbox[c]{\hsize}{\vskip7pt {\lantxh 乳杆菌属\\肠道益生菌，能够生成乳酸，抑制有害菌及炎症，调节肠道环境} \vskip7pt} & \parbox[c]{\hsize}{\vskip7pt\centerline{\raisebox{-1.5ex}{\includegraphics[width=5cm,height=0.55cm]{bar.pdf}}}\vskip7pt}  &
\hspace*{-1.51cm}\raisebox{-0.45ex}{\includegraphics{cry.pdf}}
 & \begin{minipage}{4.60cm}\begin{center}{{\color{red}\lantxh 低{\\ \bahao 不利于抑制肠道炎症及调节肠道环境}} }\end{center} \end{minipage} \\
\hline
\parbox[c]{\hsize}{\vskip7pt {\lantxh 肠球菌属\\多数为共生菌，产生乳酸。少数菌种可能导致多种感染症状} \vskip7pt} & \parbox[c]{\hsize}{\vskip7pt\centerline{\raisebox{-1.5ex}{\includegraphics[width=5cm,height=0.55cm]{bar.pdf}}}\vskip7pt}  &
\hspace*{-1.51cm}\raisebox{-0.45ex}{\includegraphics{cry.pdf}}
 & \begin{minipage}{4.60cm}\begin{center}{{\color{red}\lantxh 低{\\ \bahao 不利于肠道菌群平衡}} }\end{center} \end{minipage} \\
\hline
\parbox[c]{\hsize}{\vskip7pt {\lantxh 放线菌属\\多数为共生菌，少数在免疫力弱时可引起内源性感染、龋齿等} \vskip7pt} & \parbox[c]{\hsize}{\vskip7pt\centerline{\raisebox{-1.5ex}{\includegraphics[width=5cm,height=0.55cm]{bar.pdf}}}\vskip7pt}  &
\hspace*{-1.51cm}\raisebox{-0.45ex}{\includegraphics{cry.pdf}}
 & \begin{minipage}{4.60cm}\begin{center}{{\color{red}\lantxh 低{\\ \bahao 不利于肠道菌群平衡}} }\end{center} \end{minipage} \\
\hline
\parbox[c]{\hsize}{\vskip7pt {\lantxh 毛杆菌属\\肠道共生菌，发酵葡萄糖产生乳酸及少量乙酸和丁酸} \vskip7pt} & \parbox[c]{\hsize}{\vskip7pt\centerline{\raisebox{-1.5ex}{\includegraphics[width=5cm,height=0.55cm]{bar.pdf}}}\vskip7pt}  &
\hspace*{-1.51cm}\raisebox{-0.45ex}{\includegraphics{cry.pdf}}
 & \begin{minipage}{4.60cm}\begin{center}{{\color{red}\lantxh 低{\\ \bahao 不利于肠道菌群平衡}} }\end{center} \end{minipage} \\
\hline
\parbox[c]{\hsize}{\vskip7pt {\lantxh 芽孢杆菌属\\多数为共生菌，降解复杂碳水化合物。少数可引起感染、食物中毒等} \vskip7pt} & \parbox[c]{\hsize}{\vskip7pt\centerline{\raisebox{-1.5ex}{\includegraphics[width=5cm,height=0.55cm]{bar.pdf}}}\vskip7pt}  &
\hspace*{-1.51cm}\raisebox{-0.45ex}{\includegraphics{cry.pdf}}
 & \begin{minipage}{4.60cm}\begin{center}{{\color{red}\lantxh 低{\\ \bahao 不利于肠道菌群平衡}} }\end{center} \end{minipage} \\
\hline
\parbox[c]{\hsize}{\vskip7pt {\lantxh 阿德勒氏菌属\\帮助分解食物中的异黄酮产生雌马酚，有利于肠道及人体健康} \vskip7pt} & \parbox[c]{\hsize}{\vskip7pt\centerline{\raisebox{-1.5ex}{\includegraphics[width=5cm,height=0.55cm]{bar.pdf}}}\vskip7pt}  &
\hspace*{-1.51cm}\raisebox{-0.45ex}{\includegraphics{cry.pdf}}
 & \begin{minipage}{4.60cm}\begin{center}{{\color{red}\lantxh 低{\\ \bahao 不利于肠道及人体健康}} }\end{center} \end{minipage} \\
\hline
\parbox[c]{\hsize}{\vskip7pt {\lantxh 乳球菌属\\发酵葡萄糖产生乳酸，调节肠道pH，抑制有害菌的生长} \vskip7pt} & \parbox[c]{\hsize}{\vskip7pt\centerline{\raisebox{-1.5ex}{\includegraphics[width=5cm,height=0.55cm]{bar.pdf}}}\vskip7pt}  &
\hspace*{-1.51cm}\raisebox{-0.45ex}{\includegraphics{cry.pdf}}
 & \begin{minipage}{4.60cm}\begin{center}{{\color{red}\lantxh 低{\\ \bahao 不利于调节肠道pH及抑制有害菌生长}} }\end{center} \end{minipage} \\
\hline
\parbox[c]{\hsize}{\vskip7pt {\lantxh 真杆菌属\\产生短链脂肪酸等有益物质，少数可能与炎症、菌血症等相关} \vskip7pt} & \parbox[c]{\hsize}{\vskip7pt\centerline{\raisebox{-1.5ex}{\includegraphics[width=5cm,height=0.55cm]{bar.pdf}}}\vskip7pt}  &
\hspace*{-1.51cm}\raisebox{-0.45ex}{\includegraphics{cry.pdf}}
 & \begin{minipage}{4.60cm}\begin{center}{{\color{red}\lantxh 低{\\ \bahao 不利于产生有益物质}} }\end{center} \end{minipage} \\
\hline
\parbox[c]{\hsize}{\vskip7pt {\lantxh 埃希氏菌属\\多数为共生菌，产生维生素K，少数与腹泻、尿道感染、败血症等相关} \vskip7pt} & \parbox[c]{\hsize}{\vskip7pt\centerline{\raisebox{-1.5ex}{\includegraphics[width=5cm,height=0.55cm]{bar.pdf}}}\vskip7pt}  &
\hspace*{-1.51cm}\raisebox{-0.45ex}{\includegraphics{cry.pdf}}
 & \begin{minipage}{4.60cm}\begin{center}{{\color{red}\lantxh 低{\\ \bahao 不利于肠道菌群平衡}} }\end{center} \end{minipage} \\
\hline
\parbox[c]{\hsize}{\vskip7pt {\lantxh 别样杆菌属\\共生菌，降解多糖产生琥珀酸和乙酸，可能与长期高脂饮食相关} \vskip7pt} & \parbox[c]{\hsize}{\vskip7pt\centerline{\raisebox{-1.5ex}{\includegraphics[width=5cm,height=0.55cm]{bar.pdf}}}\vskip7pt}  &
\hspace*{-1.51cm}\raisebox{-0.45ex}{\includegraphics{cry.pdf}}
 & \begin{minipage}{4.60cm}\begin{center}{{\color{red}\lantxh 低{\\ \bahao 不利于产生有益物质}} }\end{center} \end{minipage} \\
\hline
\parbox[c]{\hsize}{\vskip7pt {\lantxh 丁酸弧菌属\\降解多种碳水化合物，产生丁酸，参与蛋白质和脂类代谢} \vskip7pt} & \parbox[c]{\hsize}{\vskip7pt\centerline{\raisebox{-1.5ex}{\includegraphics[width=5cm,height=0.55cm]{bar.pdf}}}\vskip7pt}  &
\hspace*{-1.51cm}\raisebox{-0.45ex}{\includegraphics{cry.pdf}}
 & \begin{minipage}{4.60cm}\begin{center}{{\color{red}\lantxh 低{\\ \bahao 不利于产生有益物质及物质代谢}} }\end{center} \end{minipage} \\
\hline
\end{longtable}

\vspace*{2mm}
\noindent
{\noindent\bf\xiaosihao 3、肠道致病菌}
\vspace*{1.5mm}

\begin{tctabularx}{tabularx={m{4.8cm}<{\centering}m{5.2cm}<{\centering}m{4.55cm}<{\centering}}}
&&
\\[-6pt]
\cellcolor{topcolor} \raisebox{2.614pt}{\color{black00}\bf 检测项目} &
\cellcolor{topcolor} \raisebox{2.614pt}{\color{black00}\bf 健康状态提示} &
\cellcolor{topcolor} \raisebox{2.614pt}{\color{black00}\bf 检测结果评价}
\end{tctabularx}

\vspace*{-4.25mm}
\fontsize{8.8pt}{11pt}\selectfont
\arrayrulecolor{gray2}
\begin{longtable}{m{4.8cm}m{5.2cm}<{\centering}m{0cm}@{}m{4.61cm}<{\centering}}
\hline
\parbox[c]{\hsize}{\vskip7pt {\lantxh 脆弱拟杆菌\\可能导致菌血症、腹内感染、腹膜炎} \vskip7pt} & \parbox[c]{\hsize}{\vskip7pt\centerline{\raisebox{-1.5ex}{\includegraphics[width=5cm,height=0.55cm]{bar.pdf}}}\vskip7pt}  &
\hspace*{-4.83cm}\raisebox{-0.45ex}{\includegraphics{smile.pdf}}
 & \begin{minipage}{4.60cm}\begin{center}{{\lantxh 含量未超标{}} }\end{center} \end{minipage} \\
\hline
\parbox[c]{\hsize}{\vskip7pt {\lantxh 胶胨罗斯氏菌\\可能导致败血症、脑膜炎、肺炎、菌血症等} \vskip7pt} & \parbox[c]{\hsize}{\vskip7pt\centerline{\raisebox{-1.5ex}{\includegraphics[width=5cm,height=0.55cm]{bar.pdf}}}\vskip7pt}  &
\hspace*{-4.83cm}\raisebox{-0.45ex}{\includegraphics{smile.pdf}}
 & \begin{minipage}{4.60cm}\begin{center}{{\lantxh 含量未超标{}} }\end{center} \end{minipage} \\
\hline
\parbox[c]{\hsize}{\vskip7pt {\lantxh 产气荚膜梭菌\\可导致气性坏疽，食物中毒} \vskip7pt} & \parbox[c]{\hsize}{\vskip7pt\centerline{\raisebox{-1.5ex}{\includegraphics[width=5cm,height=0.55cm]{bar.pdf}}}\vskip7pt}  &
\hspace*{-4.83cm}\raisebox{-0.45ex}{\includegraphics{smile.pdf}}
 & \begin{minipage}{4.60cm}\begin{center}{{\lantxh 含量未超标{}} }\end{center} \end{minipage} \\
\hline
\parbox[c]{\hsize}{\vskip7pt {\lantxh 病原性大肠埃希氏菌\\可能导致肠胃炎、尿路感染、新生儿脑膜炎、腹膜炎等} \vskip7pt} & \parbox[c]{\hsize}{\vskip7pt\centerline{\raisebox{-1.5ex}{\includegraphics[width=5cm,height=0.55cm]{bar.pdf}}}\vskip7pt}  &
\hspace*{-4.83cm}\raisebox{-0.45ex}{\includegraphics{smile.pdf}}
 & \begin{minipage}{4.60cm}\begin{center}{{\lantxh 含量未超标{}} }\end{center} \end{minipage} \\
\hline
\parbox[c]{\hsize}{\vskip7pt {\lantxh 副流感嗜血杆菌\\可导致脑膜炎、咽炎、心内膜炎、肺炎、败血症等} \vskip7pt} & \parbox[c]{\hsize}{\vskip7pt\centerline{\raisebox{-1.5ex}{\includegraphics[width=5cm,height=0.55cm]{bar.pdf}}}\vskip7pt}  &
\hspace*{-4.83cm}\raisebox{-0.45ex}{\includegraphics{smile.pdf}}
 & \begin{minipage}{4.60cm}\begin{center}{{\lantxh 含量未超标{}} }\end{center} \end{minipage} \\
\hline
\parbox[c]{\hsize}{\vskip7pt {\lantxh 微黄奈瑟氏球菌\\可能导致泌尿生殖系统化脓性感染、术后脑膜炎} \vskip7pt} & \parbox[c]{\hsize}{\vskip7pt\centerline{\raisebox{-1.5ex}{\includegraphics[width=5cm,height=0.55cm]{bar.pdf}}}\vskip7pt}  &
\hspace*{-4.83cm}\raisebox{-0.45ex}{\includegraphics{smile.pdf}}
 & \begin{minipage}{4.60cm}\begin{center}{{\lantxh 含量未超标{}} }\end{center} \end{minipage} \\
\hline
\parbox[c]{\hsize}{\vskip7pt {\lantxh 龋齿罗斯氏菌\\可能引起口腔疾病和其他部位的感染} \vskip7pt} & \parbox[c]{\hsize}{\vskip7pt\centerline{\raisebox{-1.5ex}{\includegraphics[width=5cm,height=0.55cm]{bar.pdf}}}\vskip7pt}  &
\hspace*{-4.83cm}\raisebox{-0.45ex}{\includegraphics{smile.pdf}}
 & \begin{minipage}{4.60cm}\begin{center}{{\lantxh 含量未超标{}} }\end{center} \end{minipage} \\
\hline
\end{longtable}

%**************************
\vspace*{2mm}
\noindent
{\noindent\bf\xiaosihao 4、肠道菌营养功能}
\vspace*{1.5mm}

\begin{tctabularx}{tabularx={m{4.8cm}<{\centering}m{5.2cm}<{\centering}m{4.55cm}<{\centering}}}
&&
\\[-6pt]
\cellcolor{topcolor} \raisebox{2.614pt}{\color{black00}\bf 检测项目} &
\cellcolor{topcolor} \raisebox{2.614pt}{\color{black00}\bf 健康状态提示} &
\cellcolor{topcolor} \raisebox{2.614pt}{\color{black00}\bf 检测结果评价}
\end{tctabularx}

\vspace*{-4.25mm}
\fontsize{8.8pt}{11pt}\selectfont
\arrayrulecolor{gray2}
\begin{longtable}{m{4.8cm}m{5.2cm}<{\centering}m{0cm}@{}m{4.61cm}<{\centering}}
\hline
\parbox[c]{\hsize}{\vskip7pt {\lantxh 维生素B1\\参与糖代谢，维持神经系统功能正常} \vskip7pt} & \parbox[c]{\hsize}{\vskip7pt\centerline{\raisebox{-1.5ex}{\includegraphics[width=5cm,height=0.55cm]{bar.pdf}}}\vskip7pt}  &
\hspace*{-1.51cm}\raisebox{-0.45ex}{\includegraphics{cry.pdf}}
 & \begin{minipage}{4.60cm}\begin{center}{{\color{red}\lantxh 低{\\ \bahao 可能引起食欲减退、乏力、头痛、肌肉酸痛}} }\end{center} \end{minipage} \\
\hline
\parbox[c]{\hsize}{\vskip7pt {\lantxh 维生素B2\\促进代谢，维护皮肤和细胞膜的完整性，抗氧化} \vskip7pt} & \parbox[c]{\hsize}{\vskip7pt\centerline{\raisebox{-1.5ex}{\includegraphics[width=5cm,height=0.55cm]{bar.pdf}}}\vskip7pt}  &
\hspace*{-1.51cm}\raisebox{-0.45ex}{\includegraphics{cry.pdf}}
 & \begin{minipage}{4.60cm}\begin{center}{{\color{red}\lantxh 低{\\ \bahao 可能增加口腔与生殖器官炎症风险}} }\end{center} \end{minipage} \\
\hline
\parbox[c]{\hsize}{\vskip7pt {\lantxh 维生素B3\\参与体内脂质代谢，降脂，扩张血管} \vskip7pt} & \parbox[c]{\hsize}{\vskip7pt\centerline{\raisebox{-1.5ex}{\includegraphics[width=5cm,height=0.55cm]{bar.pdf}}}\vskip7pt}  &
\hspace*{-1.51cm}\raisebox{-0.45ex}{\includegraphics{cry.pdf}}
 & \begin{minipage}{4.60cm}\begin{center}{{\color{red}\lantxh 低{\\ \bahao 可能增加黏膜炎症风险}} }\end{center} \end{minipage} \\
\hline
\parbox[c]{\hsize}{\vskip7pt {\lantxh 维生素B5\\参与能量代谢及抗体合成，维持皮肤及头发健康} \vskip7pt} & \parbox[c]{\hsize}{\vskip7pt\centerline{\raisebox{-1.5ex}{\includegraphics[width=5cm,height=0.55cm]{bar.pdf}}}\vskip7pt}  &
\hspace*{-1.51cm}\raisebox{-0.45ex}{\includegraphics{cry.pdf}}
 & \begin{minipage}{4.60cm}\begin{center}{{\color{red}\lantxh 低{\\ \bahao 不利于皮肤健康}} }\end{center} \end{minipage} \\
\hline
\parbox[c]{\hsize}{\vskip7pt {\lantxh 维生素B6\\抗感染，参与脂类、糖及蛋白质的代谢，参与血红蛋白合成} \vskip7pt} & \parbox[c]{\hsize}{\vskip7pt\centerline{\raisebox{-1.5ex}{\includegraphics[width=5cm,height=0.55cm]{bar.pdf}}}\vskip7pt}  &
\hspace*{-1.51cm}\raisebox{-0.45ex}{\includegraphics{cry.pdf}}
 & \begin{minipage}{4.60cm}\begin{center}{{\color{red}\lantxh 低{\\ \bahao 可能引起唇干裂、脂溢性皮炎}} }\end{center} \end{minipage} \\
\hline
\parbox[c]{\hsize}{\vskip7pt {\lantxh 维生素B7\\维持皮肤及头发健康，增强免疫，与暗视力有关} \vskip7pt} & \parbox[c]{\hsize}{\vskip7pt\centerline{\raisebox{-1.5ex}{\includegraphics[width=5cm,height=0.55cm]{bar.pdf}}}\vskip7pt}  &
\hspace*{-1.51cm}\raisebox{-0.45ex}{\includegraphics{cry.pdf}}
 & \begin{minipage}{4.60cm}\begin{center}{{\color{red}\lantxh 低{\\ \bahao 可能引发脱发和白发，不利于暗视力}} }\end{center} \end{minipage} \\
\hline
\parbox[c]{\hsize}{\vskip7pt {\lantxh 维生素B9\\参与氨基酸代谢，参与血红蛋白及核酸的合成，促进生长发育} \vskip7pt} & \parbox[c]{\hsize}{\vskip7pt\centerline{\raisebox{-1.5ex}{\includegraphics[width=5cm,height=0.55cm]{bar.pdf}}}\vskip7pt}  &
\hspace*{-1.51cm}\raisebox{-0.45ex}{\includegraphics{cry.pdf}}
 & \begin{minipage}{4.60cm}\begin{center}{{\color{red}\lantxh 低{\\ \bahao 可能导致巨幼红细胞性贫血、高同型半胱氨酸血症}} }\end{center} \end{minipage} \\
\hline
\parbox[c]{\hsize}{\vskip7pt {\lantxh 类固醇\\参与机体物质代谢，调节免疫功能} \vskip7pt} & \parbox[c]{\hsize}{\vskip7pt\centerline{\raisebox{-1.5ex}{\includegraphics[width=5cm,height=0.55cm]{bar.pdf}}}\vskip7pt}  &
\hspace*{-1.51cm}\raisebox{-0.45ex}{\includegraphics{cry.pdf}}
 & \begin{minipage}{4.60cm}\begin{center}{{\color{red}\lantxh 低{\\ \bahao 可能影响机体物质代谢，降低抵御疾病的能力}} }\end{center} \end{minipage} \\
\hline
\parbox[c]{\hsize}{\vskip7pt {\lantxh 辅酶Q\\激活细胞呼吸代谢，抗氧化，增强免疫力} \vskip7pt} & \parbox[c]{\hsize}{\vskip7pt\centerline{\raisebox{-1.5ex}{\includegraphics[width=5cm,height=0.55cm]{bar.pdf}}}\vskip7pt}  &
\hspace*{-1.51cm}\raisebox{-0.45ex}{\includegraphics{cry.pdf}}
 & \begin{minipage}{4.60cm}\begin{center}{{\color{red}\lantxh 低{\\ \bahao 可能降低免疫力及抗衰老能力}} }\end{center} \end{minipage} \\
\hline
\parbox[c]{\hsize}{\vskip7pt {\lantxh 胆汁酸\\促进食物中脂类和脂溶性维生素的吸收} \vskip7pt} & \parbox[c]{\hsize}{\vskip7pt\centerline{\raisebox{-1.5ex}{\includegraphics[width=5cm,height=0.55cm]{bar.pdf}}}\vskip7pt}  &
\hspace*{-1.51cm}\raisebox{-0.45ex}{\includegraphics{cry.pdf}}
 & \begin{minipage}{4.60cm}\begin{center}{{\color{red}\lantxh 低{\\ \bahao 不利于摄取食物中的脂类与脂溶性维生素}} }\end{center} \end{minipage} \\
\hline
\multicolumn{3}{l}{}\\
\multicolumn{3}{l}{\includegraphics[scale=1]{greencircle.pdf} {\fontsize{9.3pt}{11pt}\selectfont\bf 绿色表示健康}   \includegraphics[scale=1]{orangecircle.pdf}  {\fontsize{9.3pt}{11pt}\selectfont\bf 橙色表示需要关注}  \includegraphics[scale=1]{redcircle.pdf} {\fontsize{9.3pt}{11pt}\selectfont\bf 红色表示有风险}}
\end{longtable}

\end{document}

